\section{2001秋}
\subsection{微分積分}
\prob{
  積の関数$f(x) = g(x)\cdot h(x)$の$n$次導関数を$g(x)$や$h(x)$およびそれらの導関数を用いて表せ。
  また、下記の関数の$n$次導関数を求めよ。
  \begin{gather}
    x^{5} e^{x}
  \end{gather}
}

\prob{
  次の積分を求めよ。
  \begin{gather}
    \int \frac{1}{x^3 + 4x^2 + 5x + 2} dx
  \end{gather}
}

\newpage
\subsection{線形代数}
\prob{
  次の行列について以下の問いに答えなさい。
  \begin{gather}
    \bA = \begin{bmatrix}
      \gra & \beta & \beta \\
      \beta & \gra & \beta \\
      \beta & \beta & \gra
    \end{bmatrix}
  \end{gather}
  \begin{enumerate}[label=(\arabic*), ref=(\arabic*), itemsep=0pt]
    \item \label{2001autumn_l1_1}この行列のすべての固有値を求め、その重複度を答えなさい。
    \item \label{2001autumn_l1_2}\ref{2001autumn_l1_1}で求めた固有値に対応するすべての固有ベクトルを求めなさい。
    \item \ref{2001autumn_l1_2}で求めた固有ベクトルを用いて$\bA$を対角化しなさい。
    \item $\gra$と$\beta$の間の関係に応じて、行列$\bA$の階数を求めなさい。
  \end{enumerate}
}

\prob{
  $n$次元ベクトルの鏡映を行う$n\tm n$の行列が$\bP = \bI - \gra\bv\bv^{\rm{T}}$のように定義されている。
  ただし、$\bv$は$n$鏡映面に垂直な$n$次元ベクトル、$\bI$は$n\tm n$の単位行列、
  $\disp\gra = \frac{\raisebox{-2pt}{2}}{\bv^{\rm{T}}\bv}$である。以下の問いに答えなさい。
  \begin{enumerate}[label=(\arabic*), ref=(\arabic*), itemsep=0pt]
    \item $\bP$は直交行列であることを示しなさい。
    \item $\bP$が任意の$n$次元ベクトル$\bu$に対して
    鏡映変換$\bP\bu$を施すとき、$\bu$の鏡映面に対して垂直な成分はその方向が反転するだけであり、
    鏡映面内の成分は変化しないことを示せ。
  \end{enumerate}
}
